%%%%%%%%%%%%%%%%%%%%%%%%%%%%%%%%%%%%%%%%%%%%%%%%%%%%%%%%%%%%%%%%%%%%%%
% How to use writeLaTeX: 
%
% You edit the source code here on the left, and the preview on the
% right shows you the result within a few seconds.
%
% Bookmark this page and share the URL with your co-authors. They can
% edit at the same time!
%
% You can upload figures, bibliographies, custom classes and
% styles using the files menu.
%
%%%%%%%%%%%%%%%%%%%%%%%%%%%%%%%%%%%%%%%%%%%%%%%%%%%%%%%%%%%%%%%%%%%%%%

\documentclass[12pt]{article}

\usepackage{sbc-template}

\usepackage{graphicx,url}

%\usepackage[brazil]{babel}   
\usepackage[utf8]{inputenc}  

     
\sloppy

\title{Instructions for Authors of SBC Conferences\\ Papers and Abstracts}

\author{Luciana P. Nedel\inst{1}, Rafael H. Bordini\inst{2}, Flávio Rech
  Wagner\inst{1}, Jomi F. Hübner\inst{3} }


\address{Instituto de Informática -- Universidade Federal do Rio Grande do Sul
  (UFRGS)\\
  Caixa Postal 15.064 -- 91.501-970 -- Porto Alegre -- RS -- Brazil
\nextinstitute
  Department of Computer Science -- University of Durham\\
  Durham, U.K.
\nextinstitute
  Departamento de Sistemas e Computação\\
  Universidade Regional de Blumenal (FURB) -- Blumenau, SC -- Brazil
  \email{\{nedel,flavio\}@inf.ufrgs.br, R.Bordini@durham.ac.uk,
  jomi@inf.furb.br}
}

\begin{document} 

\maketitle

\begin{abstract}
  This meta-paper describes the style to be used in articles and short papers
  for SBC conferences. For papers in English, you should add just an abstract
  while for the papers in Portuguese, we also ask for an abstract in
  Portuguese (``resumo''). In both cases, abstracts should not have more than
  10 lines and must be in the first page of the paper.
\end{abstract}
     
\begin{resumo} 
  Este meta-artigo descreve o estilo a ser usado na confecção de artigos e
  resumos de artigos para publicação nos anais das conferências organizadas
  pela SBC. É solicitada a escrita de resumo e abstract apenas para os artigos
  escritos em português. Artigos em inglês deverão apresentar apenas abstract.
  Nos dois casos, o autor deve tomar cuidado para que o resumo (e o abstract)
  não ultrapassem 10 linhas cada, sendo que ambos devem estar na primeira
  página do artigo.
\end{resumo}


\section{Introdução}

O desenvolvimento de software envolve compreender necessidades, projetar soluções, implementá-las e mantê-las ao longo do tempo. Em sistemas que atendem diferentes perfis de usuários, essas necessidades podem envolver áreas distintas e relacionadas. Como requisitos e condições de uso mudam, o software precisa evoluir sem perder os comportamentos de que seus usuários dependem. A manutenção e a evolução também influenciam a qualidade interna do sistema \cite{alenezi2021}.

Manter e ampliar um sistema pode se tornar difícil quando responsabilidades ficam espalhadas, componentes dependem diretamente de muitos outros ou regras semelhantes são implementadas em mais de um lugar. Essas condições aumentam o esforço para localizar uma regra, avaliar o alcance de uma alteração e reutilizar partes do sistema. Estudos sobre evolução de software também analisam a relação entre mudanças e atributos internos como complexidade, acoplamento e coesão \cite{alenezi2021}. A evolução exige cuidado para que novas funcionalidades não introduzam regressões no comportamento existente.

Práticas de projeto como separação de responsabilidades, encapsulamento, coesão e controle do acoplamento ajudam a organizar o código para que mudanças fiquem mais localizadas. Padrões de projeto descrevem soluções reutilizáveis para problemas recorrentes e podem orientar a implementação de funcionalidades semelhantes \cite{alhawari2022}. Seus efeitos sobre dependências e qualidade precisam ser avaliados no contexto do sistema em que são aplicados \cite{hsu2024}.

O Conecta Obras é uma plataforma digital voltada a obras e serviços relacionados. O sistema contempla os perfis de cliente, profissional, lojista, suporte e administração. Clientes buscam soluções, profissionais oferecem serviços e lojistas disponibilizam produtos, enquanto as equipes de suporte e administração atendem usuários e gerenciam a plataforma. Assim, o projeto reúne diferentes necessidades e responsabilidades em uma mesma aplicação. Nesse contexto, a pesquisa investiga a seguinte questão: como a aplicação dos padrões State e Facade no Conecta Obras afeta a organização das responsabilidades, a duplicação, a complexidade e a testabilidade do sistema?

O objetivo geral é analisar os efeitos da aplicação de State e Facade na qualidade interna do Conecta Obras, comparando as implementações inicial e refatorada. Para isso, o trabalho busca caracterizar o sistema e seu projeto inicial; identificar problemas de acoplamento, coesão, duplicação, manutenção e extensão; aplicar os padrões selecionados; verificar, por testes automatizados, a preservação do comportamento funcional; comparar as implementações com métricas estáticas, testes e análise qualitativa; e discutir os resultados, os custos das soluções e as limitações da avaliação.




\section{Fundamentação Teórica}

\subsection{Desenvolvimento de software}

O desenvolvimento de software compreende as atividades pelas quais uma necessidade é transformada em um sistema utilizável. Requisitos, projeto, implementação, verificação, implantação e manutenção são partes desse trabalho, ainda que sua ordem e seu grau de sobreposição variem conforme o processo adotado. O produto não se reduz ao código: inclui decisões sobre comportamento, estrutura e formas de verificar se o sistema atende ao que foi solicitado. Depois da entrega, correções e novas demandas continuam a alterar o software. Por isso, o ciclo de vida abrange tanto a construção inicial quanto sua evolução \cite{hsu2024}.

Um processo de desenvolvimento define como essas atividades são organizadas, acompanhadas e revistas. Em processos tradicionais, costuma-se estabelecer um planejamento mais detalhado antes da implementação e formalizar a passagem entre etapas. Essa organização facilita o registro de decisões e a previsão de entregas quando os requisitos estão relativamente estáveis, mas torna mais custosa a incorporação de mudanças descobertas tardiamente. As abordagens ágeis trabalham com incrementos menores, avaliações frequentes do produto e ajuste das prioridades a partir do retorno obtido. A distinção não significa ausência de planejamento em um caso ou de revisões no outro; diz respeito, sobretudo, ao momento em que se assumem compromissos e à frequência com que eles são reavaliados.

A qualidade deve ser considerada durante todo o ciclo de vida. A verificação funcional examina se o sistema executa as operações esperadas; a qualidade interna diz respeito, entre outros aspectos, à facilidade de compreender, testar e modificar sua implementação. Uma funcionalidade pode estar correta hoje e, ainda assim, exigir alterações extensas para atender a um requisito novo. Estudos sobre evolução e manutenção mostram a utilidade de examinar a estrutura do código juntamente com seu comportamento, pois os efeitos de uma mudança dependem das dependências existentes e da tarefa realizada \cite{prechelt2001,hsu2024}.

\subsection{Projeto e design de software}

O design de software é a definição da estrutura de uma solução antes e durante sua implementação: quais responsabilidades pertencem a cada componente, por quais interfaces eles se comunicam e onde ficam as regras do domínio. Essas escolhas influenciam o alcance de uma alteração. Quando uma regra está distribuída por vários pontos, sua revisão exige localizar todos eles; quando as dependências são explícitas e as responsabilidades estão bem delimitadas, é mais fácil identificar o que precisa mudar \cite{gamma1994}.

Coesão descreve o quanto as responsabilidades de um componente se relacionam entre si. Acoplamento expressa o grau de dependência entre componentes. Alta coesão e acoplamento controlado são objetivos de projeto porque favorecem alterações localizadas, mas não se alcançam apenas pelo aumento do número de classes. Dividir uma operação simples entre muitos objetos pode deslocar a complexidade para as interações. Da mesma forma, concentrar operações sem relação em uma única classe dificulta compreender sua finalidade. Em uma análise empírica, o tamanho das classes mostrou influência maior sobre a propensão a mudanças do que o papel exercido por elas em padrões de projeto \cite{posnett2011}. Esse resultado recomenda cautela ao atribuir efeitos de qualidade a uma única decisão estrutural.

Encapsulamento consiste em manter detalhes e invariantes de um componente sob seu controle, expondo operações pelas quais outros componentes podem utilizá-lo. Abstração permite descrever essas operações por contratos relevantes para o cliente, sem exigir conhecimento de cada implementação concreta. A decomposição divide o problema em partes tratáveis; a separação de responsabilidades orienta essa divisão para que cada parte tenha um motivo identificável para mudar. São princípios relacionados, mas distintos: uma interface abstrata pode esconder detalhes e ainda oferecer operações demais, enquanto uma decomposição em vários módulos pode preservar dependências excessivas entre eles \cite{gamma1994}.

A relação entre design e qualidade, portanto, deve ser examinada por atributo e por contexto. Interfaces e responsabilidades bem definidas podem facilitar substituições e testes isolados; também podem acrescentar camadas de indireção que dificultam rastrear a execução. Uma análise das estruturas de interação entre classes discute riscos para a testabilidade \cite{baudry2003}; outra investiga mudanças nas dependências durante a evolução de exemplos com padrões \cite{hsu2024}. Nenhum desses indicadores, isoladamente, demonstra menor esforço de manutenção em todo sistema.

\subsection{Padrões de projeto de software}

Padrões de projeto, ou \emph{design patterns}, descrevem soluções recorrentes para problemas de design em um contexto determinado. Um padrão apresenta uma intenção, os participantes envolvidos e as consequências de organizar suas relações de certa maneira. Ele não é um trecho de código pronto: sua aplicação exige adaptar a estrutura proposta às regras e às restrições do sistema. O catálogo de Gamma e colaboradores reuniu 23 padrões orientados a objetos e consolidou um vocabulário para discutir essas decisões \cite{gamma1994}.

A noção de padrão veio da documentação de soluções recorrentes em arquitetura e foi incorporada ao desenvolvimento de software orientado a objetos. Após a publicação do catálogo, o tema passou a incluir estudos sobre identificação de padrões em código, formas de aplicá-los e seus efeitos na qualidade. Um mapeamento da pesquisa sobre padrões registra essa ampliação e mostra que estudos sobre um mesmo padrão podem chegar a conclusões diferentes quanto aos atributos avaliados \cite{ampatzoglou2013}.

Os problemas tratados pelos padrões são específicos. \emph{State}, por exemplo, organiza comportamentos que variam conforme o estado de um objeto; \emph{Facade} oferece uma interface para operações de um subsistema cujos detalhes os clientes não precisam coordenar diretamente \cite{gamma1994}. Nesses casos, a solução pode reunir regras dispersas ou reduzir o conhecimento exigido do cliente. A reutilização está principalmente na ideia de organização, e não na cópia de uma implementação. Nomear o padrão também facilita explicar por que determinadas responsabilidades foram distribuídas entre componentes.

Esses benefícios têm custos. Novas interfaces, classes e delegações aumentam o trabalho de leitura quando a variação prevista não existe ou a operação já é simples. Em um experimento controlado, soluções com padrões favoreceram parte das tarefas de manutenção, mas alternativas mais simples tiveram melhor resultado em outras \cite{prechelt2001}. Uma revisão da literatura também relata resultados divergentes entre estudos, associados a diferenças de contexto, métricas e implementação \cite{wedyan2020}. Assim, empregar um padrão como boa prática requer identificar o problema que ele resolve, explicitar seus custos e verificar seus efeitos no código concreto.

\subsection{Classificação dos padrões de projeto}

O catálogo clássico organiza os padrões pela finalidade principal em três categorias: criacionais, estruturais e comportamentais \cite{gamma1994}. A classificação ajuda a localizar soluções para problemas semelhantes, mas não determina, por si, qual padrão convém a uma implementação.

Os \emph{padrões criacionais} tratam da criação de objetos. Eles separam, em graus diferentes, o uso de um objeto da escolha de sua classe concreta e da forma de construí-lo. \emph{Factory Method} transfere a decisão de instanciação para uma operação especializada; \emph{Abstract Factory} reúne a criação de objetos relacionados. Essas alternativas são pertinentes quando a criação direta espalharia decisões sobre tipos concretos por vários clientes \cite{gamma1994}.

Os \emph{padrões estruturais} descrevem como classes e objetos se combinam para formar estruturas maiores. \emph{Composite} permite tratar objetos individuais e composições por uma interface comum; \emph{Decorator} acrescenta responsabilidades por composição. \emph{Facade} pertence a essa categoria porque organiza o acesso a um subsistema por meio de uma interface de mais alto nível. A Facade simplifica o contato dos clientes com o subsistema, mas não elimina a necessidade de distribuir adequadamente as responsabilidades entre seus componentes internos \cite{gamma1994}.

Os \emph{padrões comportamentais} tratam da distribuição de algoritmos, decisões e comunicação entre objetos. \emph{Strategy} encapsula alternativas de algoritmo; \emph{Observer} organiza a comunicação entre um objeto e seus observadores. Em \emph{State}, o contexto delega o comportamento a um objeto que representa seu estado atual, permitindo que operações e transições variem conforme esse estado \cite{gamma1994}. Essa organização pode tornar regras de estado mais localizadas, mas a quantidade de interações e a posição das transições ainda precisam ser consideradas nos testes \cite{baudry2003}.

\bibliographystyle{sbc}
\bibliography{referencias}

\end{document}
